% GENERATED FILE. Edit canonical lessons, then run npm run generate:books.
% canonical-source-hash: fnv1a64:2ffac11b1a2438f3
% canonical-lessons: GE-C452-verpackung, GE-C452-verpacken, GE-C452-karton, GE-C452-liefern, GE-C452-umtauschen

\chapter{Bestellt und geliefert}
\label{ch:ge-bestellt-und-geliefert}
\hlchaptermodality{452}

\begin{chapteropening}
\textbf{By the end of this chapter:} \emph{I can send and return things: the packaging, wrapping up, the cardboard box, delivering and exchanging.}

Complete the last lesson of chapter 452: I can send and return things: the packaging, wrapping up, the cardboard box, delivering and exchanging.
\end{chapteropening}

\section[die Verpackung]{die Verpackung — the packaging}

\label{lesson:GE-C452-verpackung}



\begin{quote}
Before the new one: say the German for to haggle, then the German for the flea market.
\end{quote}

\subsection*{You'll want to know: die Verpackung}
\textbf{die Verpackung} — \textquotedblleft{}the packaging\textquotedblright{}.

\begin{quote}
Die Verpackung ist größer als das Geschenk.
\end{quote}

\begin{quote}
\emph{The packaging is bigger than the present.}
\end{quote}

Plural: \emph{die Verpackungen}.

\begin{cousinweb}
\emph{Verpackung} is built on \emph{verpacken}, from \emph{packen} 'to pack', which was borrowed from Middle Low German \emph{packen} (EtymDB).
\end{cousinweb}

\begin{grammarlens}[title={What you've built}]
One more word for telling, explaining and describing at B1.
\end{grammarlens}

\subsection*{Your turn}
\begin{itemize}
\raggedright
  \item \emph{Say it:} \emph{die Verpackung}
  \item \emph{Say it:} \emph{die Verpackung}, once more
  \item \emph{Say it:} a sentence of your own with \emph{die Verpackung}
  \item \emph{Recall:} say the German for to haggle, then the German for the flea market, then say \emph{die Verpackung} again
\end{itemize}

\begin{tcolorbox}[breakable,colback=teal!4,colframe=teal!35!black,title={Before you move on}]
What does die Verpackung mean? (\textquotedblleft{}The packaging\textquotedblright{}.) Say it once more.
\end{tcolorbox}
\section[verpacken]{verpacken — to wrap up, to pack}

\label{lesson:GE-C452-verpacken}



\begin{quote}
Before the new one: say the German for the flea market, then the German for the packaging.
\end{quote}

\subsection*{You'll want to know: verpacken}
\textbf{verpacken} — \textquotedblleft{}to wrap up, to pack\textquotedblright{}.

\begin{quote}
Im Laden verpackt die Verkäuferin das Buch als Geschenk.
\end{quote}

\begin{quote}
\emph{In the shop the assistant wraps the book as a present.}
\end{quote}

\emph{Verpacken} — to put in packaging; \emph{einpacken} — to pack, also to gift-wrap.

\begin{cousinweb}
\emph{verpacken} is built from \emph{ver-} and \emph{packen} 'to pack', which was borrowed from Middle Low German \emph{packen} (EtymDB).
\end{cousinweb}

\begin{grammarlens}[title={What you've built}]
One more word for telling, explaining and describing at B1.
\end{grammarlens}

\subsection*{Your turn}
\begin{itemize}
\raggedright
  \item \emph{Say it:} \emph{verpacken}
  \item \emph{Say it:} \emph{verpacken}, once more
  \item \emph{Say it:} a sentence of your own with \emph{verpacken}
  \item \emph{Recall:} say the German for the flea market, then the German for the packaging, then say \emph{verpacken} again
\end{itemize}

\begin{tcolorbox}[breakable,colback=teal!4,colframe=teal!35!black,title={Before you move on}]
What does verpacken mean? (\textquotedblleft{}To wrap up, to pack\textquotedblright{}.) Say it once more.
\end{tcolorbox}
\section[der Karton]{der Karton — the cardboard box; cardboard}

\label{lesson:GE-C452-karton}



\begin{quote}
Before the new one: say the German for the packaging, then the German for to wrap up.
\end{quote}

\subsection*{You'll want to know: der Karton}
\textbf{der Karton} — \textquotedblleft{}the cardboard box; cardboard\textquotedblright{}.

\begin{quote}
Für den Umzug brauchen wir viele Kartons.
\end{quote}

\begin{quote}
\emph{We need lots of cardboard boxes for the move.}
\end{quote}

Plural: \emph{die Kartons}.

\begin{cousinweb}
\emph{Karton} was borrowed from French \emph{carton} (EtymDB), which came from Italian \emph{cartone} 'stiff paper', built on \emph{carta} 'paper', from Latin \emph{charta}.
\end{cousinweb}

\begin{grammarlens}[title={What you've built}]
One more word for telling, explaining and describing at B1.
\end{grammarlens}

\subsection*{Your turn}
\begin{itemize}
\raggedright
  \item \emph{Say it:} \emph{der Karton}
  \item \emph{Say it:} \emph{der Karton}, once more
  \item \emph{Say it:} a sentence of your own with \emph{der Karton}
  \item \emph{Recall:} say the German for the packaging, then the German for to wrap up, then say \emph{der Karton} again
\end{itemize}

\begin{tcolorbox}[breakable,colback=teal!4,colframe=teal!35!black,title={Before you move on}]
What does der Karton mean? (\textquotedblleft{}The cardboard box; cardboard\textquotedblright{}.) Say it once more.
\end{tcolorbox}
\section[liefern]{liefern — to deliver; to supply}

\label{lesson:GE-C452-liefern}



\begin{quote}
Before the new one: say the German for to wrap up, then the German for the cardboard box.
\end{quote}

\subsection*{You'll want to know: liefern}
\textbf{liefern} — \textquotedblleft{}to deliver; to supply\textquotedblright{}.

\begin{quote}
Das Geschäft liefert die Möbel am Dienstag.
\end{quote}

\begin{quote}
\emph{The shop delivers the furniture on Tuesday.}
\end{quote}

\emph{Die Lieferung} — the delivery.

\begin{cousinweb}
\emph{liefern} is inherited from Middle High German \emph{lieveren} (EtymDB), usually traced through Middle Low German to French \emph{livrer} 'to hand over', which goes back to Latin \emph{līberāre} 'to set free' (Lewis \& Short).
\end{cousinweb}

\begin{grammarlens}[title={What you've built}]
One more word for telling, explaining and describing at B1.
\end{grammarlens}

\subsection*{Your turn}
\begin{itemize}
\raggedright
  \item \emph{Say it:} \emph{liefern}
  \item \emph{Say it:} \emph{liefern}, once more
  \item \emph{Say it:} a sentence of your own with \emph{liefern}
  \item \emph{Recall:} say the German for to wrap up, then the German for the cardboard box, then say \emph{liefern} again
\end{itemize}

\begin{tcolorbox}[breakable,colback=teal!4,colframe=teal!35!black,title={Before you move on}]
What does liefern mean? (\textquotedblleft{}To deliver; to supply\textquotedblright{}.) Say it once more.
\end{tcolorbox}
\section[umtauschen]{umtauschen — to exchange, to change (in a shop)}

\label{lesson:GE-C452-umtauschen}



\begin{quote}
Before the new one: say the German for the cardboard box, then the German for to deliver.
\end{quote}

\subsection*{You'll want to know: umtauschen}
\textbf{umtauschen} — \textquotedblleft{}to exchange, to change (in a shop)\textquotedblright{}.

\begin{quote}
Die Schuhe sind zu klein, deshalb tauscht er sie um.
\end{quote}

\begin{quote}
\emph{The shoes are too small, so he exchanges them.}
\end{quote}

Separable: \emph{ich tausche um}, \emph{ich habe umgetauscht}. \emph{Der Umtausch} — the exchange.

\begin{cousinweb}
\emph{umtauschen} is built from \emph{um} 'over, round' and \emph{tauschen} 'to swap', inherited from Middle High German \emph{tiuschen} (EtymDB).
\end{cousinweb}

\begin{grammarlens}[title={What you've built}]
One more word for telling, explaining and describing at B1.
\end{grammarlens}

\subsection*{Your turn}
\begin{itemize}
\raggedright
  \item \emph{Say it:} \emph{umtauschen}
  \item \emph{Say it:} \emph{umtauschen}, once more
  \item \emph{Say it:} a sentence of your own with \emph{umtauschen}
  \item \emph{Recall:} say the German for the cardboard box, then the German for to deliver, then say \emph{umtauschen} again
\end{itemize}

\begin{tcolorbox}[breakable,colback=teal!4,colframe=teal!35!black,title={Before you move on}]
What does umtauschen mean? (\textquotedblleft{}To exchange, to change (in a shop)\textquotedblright{}.) Say it once more.
\end{tcolorbox}
